\section{Related work}
\label{sec:related}

\paragraph{Depth encoding in fluorescence microscopy.}
Depth can be encoded through detection optics or structured excitation. Light-field microscopy, DiffuserCam and Miniscope3D exploit depth-dependent detection responses \cite{broxton2013,antipa2018,yanny2020}. On the excitation side, Pascucci \etal~\cite{pascucci2019} recover volumes from lateral speckle scans through axial decorrelation, with saturation improving resolution, while Yeh \etal~\cite{yeh2019} combine translated speckles with axial scanning. Yilmaz \etal~\cite{yilmaz2015} use incident-angle scanning through a fixed scattering layer for two-dimensional fluorescence recovery; related work obtains depth cues from angular correlations \cite{takasaki2014} and tilted-pattern shifts \cite{jiang2019}. These precedents motivate our random phase-plate baseline. We optimize the phase profile of a defocused plate for depth discrimination, keeping the plate fixed throughout acquisition.

\paragraph{Learning optical encoders.}
End-to-end approaches jointly optimize optics and reconstruction networks \cite{haputhanthri2022}, including DeepSTORM3D \cite{nehme2020deepstorm3d} for emitter localization and FourierNets \cite{deb2022fouriernets} for volumetric microscopy. Kellman \etal~\cite{kellman2019} optimize illumination for a specified reconstruction algorithm. Decoder-free design instead targets measurement properties. Shechtman \etal~\cite{shechtman2014} use Fisher information to design point-spread functions (PSFs) for three-dimensional localization. Miniscope3D~\cite{yanny2020} aggregates cross-coherence between depth-dependent detection PSFs, using a smoothed maximum over lateral shifts for each depth pair. Squared-correlation objectives have a classical frame-potential interpretation \cite{benedetto2003}, and related Gram-matrix criteria are established in sensing-matrix design \cite{duarte2009}. IDEAL \cite{pinkard2025} and IDEAL-IO \cite{markley2025} optimize estimated mutual information using fitted measurement distributions, without a reconstruction decoder. Our objective specializes coherence-based design to detector-integrated voxel signatures across excitation angles, comparing depths at the same lateral position. Training uses the optical model alone, without specimen examples or a fitted measurement distribution.

\paragraph{Volumetric reconstruction and priors.}
Reconstruction methods range from total variation regularization \cite{rudin1992} to plug-and-play denoisers \cite{venkatakrishnan2013} and learned networks following a model-based inverse \cite{jin2017}. Multi-slice fusion \cite{majee2021} combines lower-dimensional denoisers along complementary orientations of a shared volume. We use these complementary approaches to assess whether gains from learned illumination persist across different reconstruction priors.
